% 实证研究 Agent 设计文档集 —— LaTeX 前导（pandoc -H）
% 用途：ctexart + xelatex，中文正文/标题/代码、彩色链接、页眉页码、目录。
\usepackage{fancyhdr}
\pagestyle{fancy}
\fancyhf{}
\fancyhead[L]{\small 实证研究 Agent · 设计文档集}
\fancyhead[R]{\small \thepage}
\renewcommand{\headrulewidth}{0.4pt}
\fancypagestyle{plain}{\fancyhf{}\fancyfoot[C]{\small \thepage}}

% 代码块字体略小以免越界；等宽用覆盖 拉丁+CJK注释+框线 的 Noto Sans SC
\usepackage{verbatim}
\makeatletter
\renewcommand{\verbatim@font}{\ttfamily\footnotesize}
\makeatother
\setmonofont{Noto Sans SC}[Scale=0.88]

% 颜色链接（避免整页蓝字，改为深蓝链接+书签）
\usepackage{xcolor}
\definecolor{linkblue}{HTML}{1F4E79}
\definecolor{softgray}{HTML}{6B7280}
\usepackage{hyperref}
\hypersetup{colorlinks=true, linkcolor=linkblue, urlcolor=linkblue,
            citecolor=linkblue, bookmarks=true, pdftitle={实证研究 Agent 设计文档集}}

% 标题小修饰
\usepackage{titlesec}
\titlespacing*{\chapter}{0pt}{0pt}{18pt}

% 显式中文字体（正文衬线 + 标题/无衬线），并把 数学符号/箭头/圆圈数字/框线
% 划到 CJK 字形范围，避免落进拉丁字体缺字形
\setCJKmainfont{Noto Serif SC}
\setCJKsansfont{Noto Sans SC}
\xeCJKDeclareCharClass{CJK}{"2190->"21FF}   % 箭头
\xeCJKDeclareCharClass{CJK}{"2200->"2211}   % 数学运算符（前半，不含减号 −）
\xeCJKDeclareCharClass{CJK}{"2213->"22FF}   % 数学运算符（后半 ≥ ≤ ≠ 等）
\xeCJKDeclareCharClass{CJK}{"2460->"24FF}   % 圆圈数字 ① ② ③
\xeCJKDeclareCharClass{CJK}{"2500->"257F}   % 框线

% emoji 状态记号（✅ ✔ 🟡）由预处理替换成下列宏，避免缺字形
\usepackage{pifont}
\definecolor{mgreen}{HTML}{1E7E34}
\definecolor{mamber}{HTML}{B45309}
\newcommand{\yesmark}{{\color{mgreen}\ding{51}}}   % ✅/✔
\newcommand{\warnmark}{{\color{mamber}\textbullet}} % 🟡
